% ==========================================================================
%  Decide, Don't Generate -- print design
%  7 x 10 in, twoside, 4.7in text block with a 1.1in outer margin column.
% ==========================================================================

% ---------- Fonts ---------------------------------------------------------
\setmainfont{Source Serif 4}[
  UprightFont = *-Regular, BoldFont = *-Semibold, ItalicFont = *-Italic,
  BoldItalicFont = *-SemiboldItalic, Extension = .ttf, Path = ./assets/fonts/,
  Numbers = {Proportional,OldStyle}, Ligatures = TeX]
\setsansfont{Inter}[
  UprightFont = *-Regular, BoldFont = *-Bold, ItalicFont = *-Italic,
  BoldItalicFont = *-BoldItalic, Extension = .ttf, Path = ./assets/fonts/,
  Numbers = Lining, Scale = 0.92]
\setmonofont{JetBrainsMono}[
  UprightFont = *-Regular, BoldFont = *-Bold, ItalicFont = *-Italic,
  BoldItalicFont = *-BoldItalic, Extension = .ttf, Path = ./assets/fonts/,
  Scale = 0.84, RawFeature = {-calt;-liga}]
\newfontfamily\semibold{Inter}[UprightFont = *-SemiBold, Extension = .ttf, Path = ./assets/fonts/, Scale = 0.92]
\newfontfamily\serifsemi{SourceSerif4}[UprightFont = *-Semibold, Extension = .ttf, Path = ./assets/fonts/, Numbers = Lining]
\newfontfamily\liningserif{SourceSerif4}[UprightFont = *-Regular, ItalicFont = *-Italic, Extension = .ttf, Path = ./assets/fonts/, Numbers = Lining]

\microtypesetup{protrusion=true,expansion=true}
\linespread{1.08}
\setlength{\parindent}{0pt}
\setlength{\parskip}{0.5em plus 0.15em minus 0.05em}
\raggedbottom
\frenchspacing
\widowpenalty=10000
\clubpenalty=10000
\brokenpenalty=5000
\emergencystretch=1.5em

% ---------- Colours (keep in sync with jevkit/figs/style.py) --------------
\definecolor{ink}{HTML}{1F2430}
\definecolor{ink2}{HTML}{4B5563}
\definecolor{muted}{HTML}{8A919C}
\definecolor{gridc}{HTML}{E4E7EB}
\definecolor{rulec}{HTML}{C9CED6}
\definecolor{datablue}{HTML}{2F6DB5}
\definecolor{llmorange}{HTML}{E07A1F}
\definecolor{jevgreen}{HTML}{3B9C6E}
\definecolor{failred}{HTML}{B8323A}
\definecolor{datatint}{HTML}{E4EDF8}
\definecolor{llmtint}{HTML}{FCEBDC}
\definecolor{jevtint}{HTML}{EEF7F2}
\definecolor{failtint}{HTML}{FBEFEF}
\definecolor{neutraltint}{HTML}{F4F5F7}
\definecolor{zoneact}{HTML}{B7A6E0}
\definecolor{zonereview}{HTML}{8468C9}
\definecolor{zoneesc}{HTML}{4B2C8F}
\definecolor{zonetint}{HTML}{F3F0FA}
\definecolor{codebg}{HTML}{F6F7F9}
\definecolor{storyyellow}{HTML}{FFF3C4}
\definecolor{storyedge}{HTML}{E0B400}
\color{ink}

% ---------- Packages --------------------------------------------------------
\usepackage{tikz}
\usetikzlibrary{calc}
\usepackage[most]{tcolorbox}
\usepackage{marginnote}
\renewcommand*{\marginfont}{\sffamily\fontsize{7}{9}\selectfont\color{ink2}}
\usepackage{changepage}
\strictpagecheck
\usepackage{enumitem}
\setlist{itemsep=2pt, topsep=4pt, parsep=0pt, leftmargin=1.4em}
\setlist[itemize]{label={\color{muted}\textbullet}}
\usepackage{needspace}
\usepackage{booktabs}
\usepackage{fvextra}

% ---------- Captions --------------------------------------------------------
\usepackage{caption}
\captionsetup{font={sf,footnotesize}, labelfont={bf,color=ink}, textfont={color=ink2},
  labelsep=period, justification=raggedright, singlelinecheck=false, skip=6pt}
\captionsetup[table]{position=top, skip=4pt}

% ---------- Page style ------------------------------------------------------
\usepackage[automark,headsepline=false]{scrlayer-scrpage}
\clearpairofpagestyles
\setkomafont{pageheadfoot}{\sffamily\fontsize{7}{9}\selectfont\color{muted}}
\setkomafont{pagenumber}{\sffamily\bfseries\fontsize{7.5}{9}\selectfont\color{ink}}
\automark[section]{chapter}
\renewcommand*{\chaptermarkformat}{\thechapter\enskip\textbar\enskip}
\renewcommand*{\sectionmarkformat}{}
\lehead{\pagemark\hspace{1.2em}\leftmark}
\rohead{\rightmark\hspace{1.2em}\pagemark}
\pagestyle{scrheadings}
\renewcommand*{\chapterpagestyle}{plain.scrheadings}
\cfoot*{}
\ofoot[\pagemark]{}

% ---------- Headings --------------------------------------------------------
\setcounter{secnumdepth}{0}
\setkomafont{disposition}{\sffamily\bfseries\color{ink}}
\setkomafont{chapter}{\sffamily\bfseries\fontsize{25}{29}\selectfont\color{ink}}
\setkomafont{section}{\sffamily\bfseries\fontsize{13.5}{17}\selectfont\color{ink}}
\setkomafont{subsection}{\sffamily\bfseries\fontsize{10.5}{13}\selectfont\color{ink}}
\setkomafont{subsubsection}{\sffamily\bfseries\fontsize{9.5}{12}\selectfont\color{ink2}}
\setkomafont{paragraph}{\sffamily\bfseries\fontsize{9.5}{12}\selectfont\color{ink}}
\RedeclareSectionCommand[beforeskip=-0.2in, afterskip=0.12in]{chapter}
\RedeclareSectionCommand[beforeskip=-1.5\baselineskip plus -4pt, afterskip=0.55\baselineskip]{section}
\RedeclareSectionCommand[beforeskip=-1.1\baselineskip plus -3pt, afterskip=0.35\baselineskip]{subsection}
\RedeclareSectionCommand[beforeskip=-0.9\baselineskip, afterskip=0.25\baselineskip]{subsubsection}
\renewcommand*{\chapterformat}{%
  {\fontsize{8.5}{10}\selectfont\color{jevgreen}\addfontfeature{LetterSpace=12}CHAPTER\ \thechapter}\par\vspace{7pt}}
\renewcommand*{\chapterlinesformat}[3]{%
  \parbox[t]{\textwidth}{\raggedright #2#3}}
\setkomafont{part}{\sffamily\bfseries\fontsize{28}{33}\selectfont\color{ink}}
\setkomafont{partnumber}{\sffamily\bfseries\fontsize{11}{13}\selectfont\color{jevgreen}}
\renewcommand*{\partformat}{\addfontfeature{LetterSpace=14}PART\ \thepart}
\renewcommand*{\raggedpart}{\raggedright}
\RedeclareSectionCommand[beforeskip=1.6in, afterskip=0.3in]{part}

% ---------- Code ------------------------------------------------------------
\newcommand{\codesize}{\fontsize{8.3}{10.8}\selectfont}
\makeatletter
\@ifundefined{Shaded}{\newenvironment{Shaded}{}{}}{}
\makeatother
\renewenvironment{Shaded}{%
  \begin{tcolorbox}[breakable, enhanced, colback=codebg, colframe=codebg, boxrule=0pt, arc=2pt,
    left=7pt, right=5pt, top=4pt, bottom=4pt, before skip=6pt, after skip=8pt,
    borderline west={1.6pt}{0pt}{rulec}]}{\end{tcolorbox}}
\RecustomVerbatimEnvironment{Highlighting}{Verbatim}{commandchars=\\\{\},fontsize=\codesize,
  breaklines, breakanywhere, breaksymbolleft={}, breaksymbolright={}}
\fvset{fontsize=\codesize, breaklines, breakanywhere, breaksymbolleft={}, breaksymbolright={}, breakindent=1.2em}
% Cell outputs (plain verbatim) get a lighter, open style
\makeatletter
\renewenvironment{verbatim}{%
  \VerbatimEnvironment
  \begin{tcolorbox}[breakable, enhanced, colback=white, colframe=gridc, boxrule=0.5pt, arc=2pt,
    left=7pt, right=5pt, top=3pt, bottom=3pt, before skip=2pt, after skip=8pt]%
  \begin{Verbatim}[fontsize=\codesize, formatcom=\color{ink2}, breaklines, breakanywhere, breaksymbolleft={}, breakindent=1.2em]}%
  {\end{Verbatim}\end{tcolorbox}}
\makeatother

% ---------- Small icons (TikZ, so they print crisp) -------------------------
\newcommand{\icontryit}[1][jevgreen]{\tikz[baseline=-0.6ex]{\draw[#1,line width=0.7pt,rounded corners=1pt] (0,-0.09) rectangle (0.3,0.13);
  \draw[#1,line width=0.7pt] (0.05,0.06)--(0.1,0.02)--(0.05,-0.02); \draw[#1,line width=0.7pt] (0.13,-0.03)--(0.21,-0.03);}}
\newcommand{\iconbreaks}[1][failred]{\tikz[baseline=-0.6ex]{\draw[#1,line width=0.7pt,line join=round] (0,-0.09)--(0.3,-0.09)--(0.15,0.16)--cycle;
  \fill[#1] (0.15,-0.045) circle (0.017); \draw[#1,line width=0.8pt] (0.15,0.0)--(0.15,0.08);}}
\newcommand{\iconthreshold}[1][zonereview]{\tikz[baseline=-0.6ex]{\draw[#1,line width=0.7pt] (0,0.02)--(0.32,0.02);
  \fill[white,draw=#1,line width=0.7pt] (0.2,0.02) circle (0.05); \draw[#1,line width=0.5pt] (0.08,-0.06)--(0.08,0.1);}}
\newcommand{\icondeeper}[1][ink2]{\tikz[baseline=-0.6ex]{\draw[#1,line width=0.6pt] (0.15,0.02) circle (0.12);
  \node[#1,font=\sffamily\bfseries\fontsize{6}{6}\selectfont] at (0.15,0.02) {$\Sigma$};}}
\newcommand{\iconkey}[1][jevgreen]{\tikz[baseline=-0.6ex]{\fill[#1] (0.1,0.02) circle (0.07);}}
\newcommand{\iconstory}{\tikz[baseline=-0.6ex]{\draw[storyedge,line width=0.7pt] (0,-0.08)--(0.24,0.14); \fill[storyedge] (0,-0.08)--(0.05,-0.06)--(0.02,-0.03)--cycle;}}

\newcommand{\boxlabel}[3]{{\sffamily\bfseries\fontsize{7}{9}\selectfont\color{#1}#2\hspace{0.45em}\addfontfeature{LetterSpace=10}#3}\par\vspace{3pt}}

% ---------- The book's boxes ------------------------------------------------
% Short boxes stay whole on one page; a split one loses its label on the second page.
\newtcolorbox{tryit}[1][]{enhanced, colback=jevtint, colframe=jevtint, boxrule=0pt, arc=2pt,
  borderline west={2.2pt}{0pt}{jevgreen}, left=9pt, right=8pt, top=7pt, bottom=7pt, before skip=10pt, after skip=10pt,
  fontupper=\small, before upper={\boxlabel{jevgreen}{\icontryit}{TRY IT}}, #1}
\newtcolorbox{breaks}[1][]{enhanced, colback=failtint, colframe=failtint, boxrule=0pt, arc=2pt,
  borderline west={2.2pt}{0pt}{failred}, left=9pt, right=8pt, top=7pt, bottom=7pt, before skip=10pt, after skip=10pt,
  fontupper=\small, before upper={\boxlabel{failred}{\iconbreaks}{WHERE THIS BREAKS}}, #1}
\newtcolorbox{threshold}[1][]{enhanced, colback=zonetint, colframe=zonetint, boxrule=0pt, arc=2pt,
  borderline west={2.2pt}{0pt}{zonereview}, left=9pt, right=8pt, top=7pt, bottom=7pt, before skip=10pt, after skip=10pt,
  fontupper=\small, before upper={\boxlabel{zonereview}{\iconthreshold}{SET THE THRESHOLD}}, #1}
\newtcolorbox{deeper}[1][]{breakable, enhanced, colback=neutraltint, colframe=neutraltint, boxrule=0pt, arc=2pt,
  borderline west={2.2pt}{0pt}{rulec}, left=9pt, right=8pt, top=7pt, bottom=7pt, before skip=10pt, after skip=10pt,
  fontupper=\small, before upper={\boxlabel{ink2}{\icondeeper}{GOING DEEPER\ \ \textmd{\textcolor{muted}{optional}}}}, #1}
\newtcolorbox{sidebar}[2][]{breakable, enhanced, colback=neutraltint, colframe=neutraltint, boxrule=0pt, arc=2pt,
  left=9pt, right=8pt, top=7pt, bottom=7pt, before skip=10pt, after skip=10pt,
  fontupper=\small, before upper={\boxlabel{ink2}{}{#2}}, #1}
\newtcolorbox{storybox}[1][]{enhanced, colback=storyyellow, colframe=storyedge, boxrule=0.6pt, arc=2pt,
  left=7pt, right=7pt, top=4pt, bottom=4pt, before skip=8pt, after skip=8pt,
  fontupper=\sffamily\footnotesize, #1}

% Key idea: a reveal line with a green rule
\newtcolorbox{keyidea}[1][]{enhanced, colback=white, colframe=white, boxrule=0pt, arc=0pt,
  borderline west={3pt}{0pt}{jevgreen}, left=10pt, right=0pt, top=2pt, bottom=2pt, before skip=10pt, after skip=10pt,
  fontupper=\sffamily\bfseries\fontsize{11}{15}\selectfont\color{ink}\raggedright, #1}

% Pull quote
\newenvironment{bookquote}[1]{\par\vspace{10pt}\needspace{4\baselineskip}\noindent
  \begin{minipage}{\linewidth}\raggedright
  {\fontsize{30}{30}\selectfont\color{jevgreen}\serifsemi\char"201C}\par\vspace{-14pt}
  \leftskip=1.2em\rightskip=1.5em plus 1fil
  \itshape\fontsize{11.5}{16}\selectfont\color{ink}\def\quoteattrib{#1}}%
  {\par\vspace{5pt}\upshape\sffamily\fontsize{7.5}{9}\selectfont\color{ink2}\textemdash\ \quoteattrib\end{minipage}\par\vspace{10pt}}

% Exercises
% The hook to the next chapter must never sit alone on a page: no page break after the exercise list either.
\makeatletter\newcommand*{\keeplistend}{\@endparpenalty=\@M}\makeatother
\newenvironment{exercises}{\par\vspace{8pt}\needspace{6\baselineskip}\keeplistend\noindent
  {\sffamily\bfseries\fontsize{7}{9}\selectfont\color{ink}\addfontfeature{LetterSpace=10}EXERCISES}\par
  \vspace{2pt}{\color{rulec}\rule{\linewidth}{0.5pt}}\par\vspace{2pt}\small}{\par\nopagebreak[4]\vspace{6pt}}

% Chapter goals ("By the end...")
\newenvironment{goals}{\par\vspace{6pt}\begin{tcolorbox}[enhanced, colback=white, colframe=gridc, boxrule=0.6pt, arc=2pt,
  left=9pt, right=8pt, top=7pt, bottom=6pt, fontupper=\small]
  {\sffamily\bfseries\fontsize{7}{9}\selectfont\color{ink}\addfontfeature{LetterSpace=10}BY THE END OF THIS CHAPTER}\par\vspace{3pt}}%
  {\end{tcolorbox}\vspace{2pt}}

% Next-chapter hook
% The hook never starts a page on its own: no break is allowed between it and what comes before.
\newenvironment{nexthook}{\par\nopagebreak[4]\vspace{9pt}\nopagebreak[4]\noindent\begin{minipage}{\linewidth}\raggedright
  {\color{rulec}\rule{1.2in}{0.5pt}}\par\vspace{6pt}\sffamily\fontsize{9.5}{13}\selectfont\color{ink}}{\end{minipage}\par}

% Wide material: extends into the outer margin on both odd and even pages
\makeatletter
% Inside a wide block, figures stop floating so they take the widened line width.
% (The print layout has no margin column, so wide figures stay inside the text block.)
\newenvironment{widefig}{\par\begin{adjustwidth*}{0pt}{0pt}%
  \renewenvironment{figure}[1][]{\par\medskip\def\@captype{figure}\centering}{\par\medskip}}%
  {\end{adjustwidth*}\par}
\makeatother

% Full-page visual summary
\newenvironment{summarypage}{\clearpage\thispagestyle{plain.scrheadings}\begin{adjustwidth*}{0pt}{-\dimexpr\marginparwidth+\marginparsep\relax}\centering}{\end{adjustwidth*}\clearpage}

% Margin items
\newcommand{\sidenote}[1]{\marginnote{\raggedright\sffamily\fontsize{7}{9.2}\selectfont\color{ink2}#1}}
% A short aside set in the text column (the print layout has no margin notes)
\newenvironment{asidenote}{\par\vspace{2pt}\begingroup\leftskip=1.2em\sffamily\fontsize{8.5}{11.5}\selectfont\color{ink2}}{\par\endgroup\vspace{2pt}}

% Draft markers (visible on purpose; removed before print)
\newcommand{\authorstory}[1]{\begin{storybox}\iconstory\hspace{0.4em}\textbf{AUTHOR STORY:} #1\end{storybox}}
\newcommand{\verifymark}{\,{\sffamily\bfseries\fontsize{6}{7}\selectfont\colorbox{storyyellow}{\color{ink}VERIFY}}}
\newcommand{\syntheticmark}{{\sffamily\bfseries\fontsize{6}{7}\selectfont\color{muted}SYNTHETIC}}

% Tables
\renewcommand{\arraystretch}{1.18}
\AtBeginEnvironment{longtable}{\sffamily\fontsize{7.6}{10}\selectfont}
\AtBeginEnvironment{tabular}{\sffamily\fontsize{7.6}{10}\selectfont}

% Footnotes
\deffootnote[1.2em]{1.2em}{1em}{\textsuperscript{\thefootnotemark}\,}
\addtokomafont{footnote}{\fontsize{8}{10.5}\selectfont}

\makeatletter
\newcommand{\noindentnext}{\@afterindentfalse\@afterheading}
\makeatother

% Part pages: keep the part's short introduction on the same page as its title
\renewcommand*{\partheadendvskip}{\vspace{0.35in}}
\renewcommand*{\partheademptypage}{}
\newenvironment{partintro}{\begin{minipage}{0.92\textwidth}\raggedright\fontsize{11.5}{17}\selectfont\color{ink2}}{\end{minipage}\cleardoubleoddemptypage}  % the part's first chapter opens on a right-hand page

% Index of terms: two columns, small type
\usepackage{multicol}
\newenvironment{indexlist}{\par\begin{multicols}{2}\small\raggedright\setlength{\parskip}{2pt}}{\end{multicols}\par}

% Chapter epigraph: a short quotation set right under the chapter title
\newenvironment{chapterepigraph}[1]{\par\vspace{-2pt}\begin{flushright}\begin{minipage}{0.72\linewidth}\raggedleft
  \itshape\fontsize{9.8}{13.5}\selectfont\color{ink2}\def\epiattrib{#1}}%
  {\par\vspace{3pt}\upshape\sffamily\fontsize{7.2}{9}\selectfont\color{muted}\textemdash\ \epiattrib\end{minipage}\end{flushright}\par\vspace{8pt}}

% Chapters and parts may open on either page: no blank pages before them (print length matters).
\KOMAoptions{open=any}
\makeatletter\@openrightfalse\makeatother
% Quarto puts a \cleardoublepage before the appendix; with open=any it should only start a new page.
\let\cleardoublepage\clearpage
% Part entries in the table of contents are bold; give their page numbers room so three digits stay inside the
% margin. (Before the patch below: redeclaring \part afterwards would drop it.)
\RedeclareSectionCommand[tocpagenumberwidth=2.4em]{part}
% Part pages open on the right, like a printed book's.
\pretocmd{\part}{\cleardoubleoddemptypage}{}{}
% Roman numerals run through the front matter and Part I's opener; Chapter 1 is page 1. Quarto calls \mainmatter
% straight after the table of contents, so that call does nothing and parts/p1.qmd starts the main matter instead.
\let\bookmainmatter\mainmatter
\renewcommand*{\mainmatter}{}

% Real space characters in the PDF, so copied or extracted text keeps every space between words.
\usepackage{tagpdf}
\tagpdfsetup{interwordspace=true}

% ---------- Index ----------
% filters/book.lua writes \index entries (terms listed by tools/make_index.py); Quarto runs makeindex with
% assets/latex/book.ist. The index page's own heading comes from back/index-terms.qmd, so theindex adds none.
\usepackage{makeidx}
\makeindex
\newcommand*{\indexletter}[1]{\par\vspace{7pt plus 2pt}{\sffamily\bfseries\small\color{ink}#1}\par\nopagebreak\vspace{1pt}}
\makeatletter
\renewenvironment{theindex}{%
  \begin{multicols}{2}\small\raggedright\setlength{\parindent}{0pt}\setlength{\parskip}{0pt plus 0.3pt}%
  \def\@idxitem{\par\hangindent 1.2em}\let\item\@idxitem
  \def\subitem{\@idxitem\hspace*{0.8em}}\def\indexspace{\par}}%
  {\end{multicols}}
\makeatother

% Reference matter in small type: the bibliography here; the glossary and cheat sheets set it in their own files.
\AtBeginEnvironment{CSLReferences}{\small}

% PDF bookmarks carry part and chapter numbers.
\AtBeginDocument{\bookmarksetup{numbered}}
